\documentclass[10pt,a4paper]{amsart}
\usepackage{amsmath,amssymb,amsthm,graphicx}
% Retained preamble for the Stacks Project chapter {CH} (shared project preamble).
% The chapter file pinned-original-source.tex only carries \input{{preamble}} on its first line,
% so no preamble text lives inside the chapter itself.  The Stacks Project uses one single shared
% preamble for every chapter: the snapshot stages it three times (sources/stacks-chow/preamble.tex,
% sources/stacks-etale-cohomology/preamble.tex and focused-seeds/AGP/sources/stacks-snapshot-20261011/
% preamble.tex) and those copies are identical in size.  The block below is copied mechanically,
% line for line, from sources/stacks-chow/preamble.tex lines 14-18 and 154-217: the Xy-pic setup and
% the native theorem declarations together with the mathematical macros.  Nothing was added, removed
% or redefined; the wrapper class and the amsmath/amssymb/amsthm/graphicx base packages are supplied
% by the assembler.
% For commutative diagrams we use Xy-pic
% Chapter: stacks-spaces-more-morphisms.tex -- More on Morphisms of Spaces.
\usepackage[all]{xy}

% We use 2cell for 2-commutative diagrams.
\xyoption{2cell}

% Theorem environments.
%
\theoremstyle{plain}
\newtheorem{theorem}[subsection]{Theorem}
\newtheorem{proposition}[subsection]{Proposition}
\newtheorem{lemma}[subsection]{Lemma}

\theoremstyle{definition}
\newtheorem{definition}[subsection]{Definition}
\newtheorem{example}[subsection]{Example}
\newtheorem{exercise}[subsection]{Exercise}
\newtheorem{situation}[subsection]{Situation}

\theoremstyle{remark}
\newtheorem{remark}[subsection]{Remark}
\newtheorem{remarks}[subsection]{Remarks}

\numberwithin{equation}{subsection}

% Macros
%
\def\lim{\mathop{\mathrm{lim}}\nolimits}
\def\colim{\mathop{\mathrm{colim}}\nolimits}
\def\Spec{\mathop{\mathrm{Spec}}}
\def\Hom{\mathop{\mathrm{Hom}}\nolimits}
\def\Ext{\mathop{\mathrm{Ext}}\nolimits}
\def\SheafHom{\mathop{\mathcal{H}\!\mathit{om}}\nolimits}
\def\SheafExt{\mathop{\mathcal{E}\!\mathit{xt}}\nolimits}
\def\Sch{\mathit{Sch}}
\def\Mor{\mathop{\mathrm{Mor}}\nolimits}
\def\Ob{\mathop{\mathrm{Ob}}\nolimits}
\def\Sh{\mathop{\mathit{Sh}}\nolimits}
\def\NL{\mathop{N\!L}\nolimits}
\def\CH{\mathop{\mathrm{CH}}\nolimits}
\def\proetale{{pro\text{-}\acute{e}tale}}
\def\etale{{\acute{e}tale}}
\def\QCoh{\mathit{QCoh}}
\def\Ker{\mathop{\mathrm{Ker}}}
\def\Im{\mathop{\mathrm{Im}}}
\def\Coker{\mathop{\mathrm{Coker}}}
\def\Coim{\mathop{\mathrm{Coim}}}

% Boxtimes
%
\DeclareMathSymbol{\boxtimes}{\mathbin}{AMSa}{"02}

%
% Macros for moduli stacks/spaces
%
\def\QCohstack{\mathcal{QC}\!\mathit{oh}}
\def\Cohstack{\mathcal{C}\!\mathit{oh}}
\def\Spacesstack{\mathcal{S}\!\mathit{paces}}
\def\Quotfunctor{\mathrm{Quot}}
\def\Hilbfunctor{\mathrm{Hilb}}
\def\Curvesstack{\mathcal{C}\!\mathit{urves}}
\def\Polarizedstack{\mathcal{P}\!\mathit{olarized}}
\def\Complexesstack{\mathcal{C}\!\mathit{omplexes}}
% \Pic is the operator that assigns to X its picard group, usage \Pic(X)
% \Picardstack_{X/B} denotes the Picard stack of X over B
% \Picardfunctor_{X/B} denotes the Picard functor of X over B
\def\Pic{\mathop{\mathrm{Pic}}\nolimits}
\def\Picardstack{\mathcal{P}\!\mathit{ic}}
\def\Picardfunctor{\mathrm{Pic}}
\def\Deformationcategory{\mathcal{D}\!\mathit{ef}}
\newenvironment{reference}{}{}
\newenvironment{slogan}{}{}
\newenvironment{history}{}{}
\usepackage{cleveref}
\crefname{originallemma}{original lemma}{original lemmas}
\Crefname{originallemma}{Original lemma}{Original lemmas}
\makeatletter
\newlabel{spaces-morphisms-definition-normalization-X-in-Y}{{0822}{source}{}{sourcepointer-1}{}}
\newlabel{spaces-morphisms-definition-normalization-X-in-Y@cref}{{[originallemma][0][]0822}{[1][0][]source}{}{}{}}
\makeatother
\crefname{originallemma}{original lemma}{original lemmas}
\Crefname{originallemma}{Original lemma}{Original lemmas}
\makeatletter
\newlabel{spaces-morphisms-lemma-compose-after-separated}{{03KR}{source}{}{sourcepointer-2}{}}
\newlabel{spaces-morphisms-lemma-compose-after-separated@cref}{{[originallemma][0][]03KR}{[1][0][]source}{}{}{}}
\makeatother
\crefname{originallemma}{original lemma}{original lemmas}
\Crefname{originallemma}{Original lemma}{Original lemmas}
\makeatletter
\newlabel{spaces-morphisms-lemma-normalization-in-universally-closed}{{0A0Q}{source}{}{sourcepointer-3}{}}
\newlabel{spaces-morphisms-lemma-normalization-in-universally-closed@cref}{{[originallemma][0][]0A0Q}{[1][0][]source}{}{}{}}
\makeatother
\crefname{originallemma}{original lemma}{original lemmas}
\Crefname{originallemma}{Original lemma}{Original lemmas}
\makeatletter
\newlabel{spaces-morphisms-lemma-pushforward}{{03M9}{source}{}{sourcepointer-4}{}}
\newlabel{spaces-morphisms-lemma-pushforward@cref}{{[originallemma][0][]03M9}{[1][0][]source}{}{}{}}
\makeatother
\crefname{originallemma}{original lemma}{original lemmas}
\Crefname{originallemma}{Original lemma}{Original lemmas}
\makeatletter
\newlabel{spaces-morphisms-lemma-quasi-compact-permanence}{{03KS}{source}{}{sourcepointer-5}{}}
\newlabel{spaces-morphisms-lemma-quasi-compact-permanence@cref}{{[originallemma][0][]03KS}{[1][0][]source}{}{}{}}
\makeatother
\crefname{originallemma}{original lemma}{original lemmas}
\Crefname{originallemma}{Original lemma}{Original lemmas}
\makeatletter
\newlabel{spaces-morphisms-lemma-universally-closed-permanence}{{04NX}{source}{}{sourcepointer-6}{}}
\newlabel{spaces-morphisms-lemma-universally-closed-permanence@cref}{{[originallemma][0][]04NX}{[1][0][]source}{}{}{}}
\makeatother
\crefname{originallemma}{original lemma}{original lemmas}
\Crefname{originallemma}{Original lemma}{Original lemmas}
\makeatletter
\newlabel{spaces-morphisms-lemma-universally-closed-quasi-compact}{{04XW}{source}{}{sourcepointer-7}{}}
\newlabel{spaces-morphisms-lemma-universally-closed-quasi-compact@cref}{{[originallemma][0][]04XW}{[1][0][]source}{}{}{}}
\makeatother
\title{The normalization factorization of a universally closed quasi-separated map of spaces, flat-base-change compatible}
\author{The Stacks Project Authors}
\date{stacks-project:spaces-more-morphisms.tex:local-snapshot-20261011; GNU FDL 1.2 or later}
\begin{document}
\maketitle
\noindent\textit{Editorial note.} Statement span is lines 7993-8017 of sources/stacks-spaces-more-morphisms/pinned-original-source.tex; proof span is lines 8019-8058, the complete author proof (38 body lines).\par\medskip
\noindent\textit{Editorial note.} Navigation stubs: 7 cited-but-not-included labels are supplied as non-load-bearing pointers with their pinned tags where the snapshot tag index has them.\par\medskip
\noindent\textit{Editorial note.} extra\_packages is empty on purpose: the retained preamble already loads xy with options and a second declaration clashes in the assembler wrapper.\par\medskip
\noindent\textit{Editorial note.} Text is retained verbatim, including the author's own typos and placeholders.\par\medskip
\noindent\textit{Editorial note.} Counted B side (both complete). Pillar consumed by both Stein theorems L16200/L16203 and the lsc card L16212 — all one wave, chain closes.\par\medskip
\noindent\textit{Allowed standard prerequisites (editorial).}\begin{itemize}
\item Explicit prerequisite boundary. Morphisms of Spaces: universally-closed-quasi-compact, normalization-in-universally-closed, universally-closed-permanence, quasi-compact-permanence, compose-after-separated, pushforward; Morphisms of Spaces definition normalization-of-Y-in-X.
\end{itemize}\medskip
\ifdefined\hypertarget\hypertarget{sourcepointer-1}{}\fi
\noindent\textit{Source reference.} Original lemma 0822: sources/stacks-spaces-more-morphisms/pinned-original-source.tex, source lines 8058--8058. Cited by the retained statement or proof but not part of the collected spans; supplied only as a navigation pointer. The argument does not depend on its content.\par\medskip
\ifdefined\hypertarget\hypertarget{sourcepointer-2}{}\fi
\noindent\textit{Source reference.} Original lemma 03KR: sources/stacks-spaces-more-morphisms/pinned-original-source.tex, source lines 8058--8058. Cited by the retained statement or proof but not part of the collected spans; supplied only as a navigation pointer. The argument does not depend on its content.\par\medskip
\ifdefined\hypertarget\hypertarget{sourcepointer-3}{}\fi
\noindent\textit{Source reference.} Original lemma 0A0Q: sources/stacks-spaces-more-morphisms/pinned-original-source.tex, source lines 8058--8058. Cited by the retained statement or proof but not part of the collected spans; supplied only as a navigation pointer. The argument does not depend on its content.\par\medskip
\ifdefined\hypertarget\hypertarget{sourcepointer-4}{}\fi
\noindent\textit{Source reference.} Original lemma 03M9: sources/stacks-spaces-more-morphisms/pinned-original-source.tex, source lines 8058--8058. Cited by the retained statement or proof but not part of the collected spans; supplied only as a navigation pointer. The argument does not depend on its content.\par\medskip
\ifdefined\hypertarget\hypertarget{sourcepointer-5}{}\fi
\noindent\textit{Source reference.} Original lemma 03KS: sources/stacks-spaces-more-morphisms/pinned-original-source.tex, source lines 8058--8058. Cited by the retained statement or proof but not part of the collected spans; supplied only as a navigation pointer. The argument does not depend on its content.\par\medskip
\ifdefined\hypertarget\hypertarget{sourcepointer-6}{}\fi
\noindent\textit{Source reference.} Original lemma 04NX: sources/stacks-spaces-more-morphisms/pinned-original-source.tex, source lines 8058--8058. Cited by the retained statement or proof but not part of the collected spans; supplied only as a navigation pointer. The argument does not depend on its content.\par\medskip
\ifdefined\hypertarget\hypertarget{sourcepointer-7}{}\fi
\noindent\textit{Source reference.} Original lemma 04XW: sources/stacks-spaces-more-morphisms/pinned-original-source.tex, source lines 8058--8058. Cited by the retained statement or proof but not part of the collected spans; supplied only as a navigation pointer. The argument does not depend on its content.\par\medskip
\begin{lemma}
\label{lemma-stein-universally-closed}
Let $S$ be a scheme. Let $f : X \to Y$ be a universally closed and
quasi-separated morphism of algebraic spaces over $S$.
There exists a factorization
$$
\xymatrix{
X \ar[rr]_{f'} \ar[rd]_f & & Y' \ar[dl]^\pi \\
& Y &
}
$$
with the following properties:
\begin{enumerate}
\item the morphism $f'$ is universally closed, quasi-compact, quasi-separated,
and surjective,
\item the morphism $\pi : Y' \to Y$ is integral,
\item we have $f'_*\mathcal{O}_X = \mathcal{O}_{Y'}$,
\item we have $Y' = \underline{\Spec}_Y(f_*\mathcal{O}_X)$, and
\item $Y'$ is the normalization of $Y$ in $X$ as defined in
Morphisms of Spaces, Definition
\ref{spaces-morphisms-definition-normalization-X-in-Y}.
\end{enumerate}
Formation of the factorization $f = \pi \circ f'$ commutes with flat
base change.
\end{lemma}

\begin{proof}
By Morphisms of Spaces, Lemma
\ref{spaces-morphisms-lemma-universally-closed-quasi-compact}
the morphism $f$ is quasi-compact.
We just define $Y'$ as the normalization of $Y$ in $X$, so (5) and (2) hold
automatically. By
Morphisms of Spaces, Lemma
\ref{spaces-morphisms-lemma-normalization-in-universally-closed}
we see that (4) holds. The morphism $f'$ is universally closed by
Morphisms of Spaces, Lemma
\ref{spaces-morphisms-lemma-universally-closed-permanence}.
It is quasi-compact by
Morphisms of Spaces, Lemma
\ref{spaces-morphisms-lemma-quasi-compact-permanence}
and quasi-separated by
Morphisms of Spaces, Lemma
\ref{spaces-morphisms-lemma-compose-after-separated}.

\medskip\noindent
To show the remaining statements we may assume the base $Y$ is affine
(as taking normalization commutes with \'etale localization).
Say $Y = \Spec(R)$. Then $Y' = \Spec(A)$ with
$A = \Gamma(X, \mathcal{O}_X)$ an integral $R$-algebra.
Thus it is clear that $f'_*\mathcal{O}_X$
is $\mathcal{O}_{Y'}$ (because $f'_*\mathcal{O}_X$ is quasi-coherent,
by Morphisms of Spaces, Lemma
\ref{spaces-morphisms-lemma-pushforward},
and hence equal to $\widetilde{A}$). This proves (3).

\medskip\noindent
Let us show that $f'$ is surjective. As $f'$ is universally closed (see above)
the image of $f'$ is a closed subset
$V(I) \subset Y' = \Spec(A)$. Pick $h \in I$. Then
$h|_X = f^\sharp(h)$ is a global section of the structure sheaf of
$X$ which vanishes at every point. As $X$ is quasi-compact this means
that $h|_X$ is a nilpotent section, i.e., $h^n|X = 0$ for some $n > 0$.
But $A = \Gamma(X, \mathcal{O}_X)$, hence $h^n = 0$.
In other words $I$ is contained in the Jacobson radical of $A$ and we conclude
that $V(I) = Y'$ as desired.
\end{proof}

\end{document}
